\begin{definition}[The full accessible subspace of a charge label]
    \label{def:peps_regular_charge_subspace}
    \lean{TNLean.PEPS.regularChargeVector,TNLean.PEPS.regularChargeSubspace}
    \leanok
    \uses{def:peps_regular_charge_coordinates}
    Regard the actual matrix $A_\chi(p;x,y)$ as a vector in
    $\ell^2(G\times G)$, with its ordinary Hilbert--Schmidt inner product.
    For a coefficient function $\chi:G\to\mathbb C$, define
    \[
        \mathcal H_\chi
        =\operatorname{span}\{A_\chi(p;x,y):p,x,y\in G\}.
    \]
    Every internal parameter and both unknown vertex translations are included.
    This is an auxiliary accessible-coordinate subspace for
    \cite[charge detection]{Schuch2010PEPS}, lines 2470--2486;
    no original-spin detector is asserted.
\end{definition}

\begin{definition}[Independent translations and the accessible charge projection]
    \label{def:peps_regular_charge_translation_detector}
    \lean{TNLean.PEPS.regularChargeTranslation,TNLean.PEPS.regularChargeDetector}
    \leanok
    \uses{def:peps_regular_charge_subspace}
    Define the unitary register permutation
    $T_{z,w}f(r,s)=f(z^{-1}r,w^{-1}s)$.
    Let $P_\chi$ be the orthogonal projection onto $\mathcal H_\chi$.
\end{definition}

\begin{theorem}[The actual vectors and their Hilbert-space pairing]
    \label{thm:peps_regular_charge_subspace_vectors}
    \lean{TNLean.PEPS.regularChargeVector_mem_subspace,
        TNLean.PEPS.regularChargeVector_inner}
    \leanok
    \uses{def:peps_regular_charge_subspace}
    Every $A_\chi(p;x,y)$ belongs to $\mathcal H_\chi$, and its
    Hilbert-space inner product with $A_\psi(q;x',y')$ is exactly
    \[
        \sum_{r,s\in G}\overline{A_\chi(p;x,y)_{r,s}}
        A_\psi(q;x',y')_{r,s}.
    \]
\end{theorem}
\begin{proof}\leanok
    Membership follows from the defining span. The inner product is the
    ordinary sum of the entrywise pairings.
\end{proof}

\begin{theorem}[The unknown translations preserve the complete charge subspace]
    \label{thm:peps_regular_charge_subspace_translation}
    \lean{TNLean.PEPS.regularChargeTranslation_vector,
        TNLean.PEPS.regularChargeSubspace_map_translation}
    \leanok
    \uses{def:peps_regular_charge_translation_detector,
        thm:peps_regular_charge_coordinates_covariance}
    For all $p,x,y,z,w\in G$,
    \[
        T_{z,w}A_\chi(p;x,y)=A_\chi(p;zx,wy),
        \qquad T_{z,w}\mathcal H_\chi=\mathcal H_\chi.
    \]
\end{theorem}
\begin{proof}\leanok
    Apply the actual row and column covariance formula. The inverse
    translations give the reverse inclusion of the subspaces.
\end{proof}

\begin{theorem}[The accessible projection and its selected-charge action]
    \label{thm:peps_regular_charge_detector}
    \lean{TNLean.PEPS.regularChargeDetector_isSymmetricProjection,
        TNLean.PEPS.regularChargeDetector_vector,
        TNLean.PEPS.regularChargeDetector_commute_translation}
    \leanok
    \uses{thm:peps_regular_charge_subspace_translation,
        thm:peps_regular_charge_subspace_vectors}
    The operator $P_\chi$ is an orthogonal projection. For every $p,x,y$,
    \[
        P_\chi A_\chi(p;x,y)=A_\chi(p;x,y),
        \qquad P_\chi T_{z,w}=T_{z,w}P_\chi.
    \]
    Thus the same accessible projection works for all unknown translations.
\end{theorem}
\begin{proof}\leanok
    Orthogonal projection fixes its range. A unitary preserving the subspace
    also preserves its orthogonal complement, and hence commutes with the
    corresponding projection.
\end{proof}

\begin{theorem}[Orthogonality of distinct full charge subspaces]
    \label{thm:peps_regular_charge_subspace_orthogonality}
    \lean{TNLean.PEPS.regularChargeSubspace_isOrtho,
        TNLean.PEPS.regularChargeDetector_other_vector}
    \leanok
    \uses{thm:peps_regular_charge_coordinates_orthogonality,
        thm:peps_regular_charge_detector}
    Let $\sigma,\tau$ be finite-dimensional irreducible complex
    representations, with $\sigma$ unitary and distinct characters.
    Then $\mathcal H_{\chi_\sigma}\perp\mathcal H_{\chi_\tau}$, and
    \[
        P_{\chi_\sigma}A_{\chi_\tau}(p;x,y)=0
        \quad\text{for every }p,x,y\in G.
    \]
    These are auxiliary charge-space identities from
    \cite[charge detection]{Schuch2010PEPS}, lines 2474--2486;
    neither completeness of a physical measurement nor its original-spin
    implementation is asserted here.
\end{theorem}
\begin{proof}\leanok
    Every generator of one span is orthogonal to every generator of the
    other by the actual charge-matrix pairing. Extend by linearity and
    apply orthogonal projection onto the first span.
\end{proof}
